\begin{lemma}
  Let $X:\CHaus$ and $C\subseteq X$ closed. 
  Then there exists some sequence of open subsets $U_n:X \to \Open$ such that $C = \bigcap_{n:\N} U_n$. 
\end{lemma}
\begin{proof}
  Assume a cover $q:S\twoheadrightarrow X$ with $S:\Stone$. Then $q^{-1}(C)\subseteq S$ is closed and by \parencite[TODO]{TODO}, equal to $\bigcap_{n:\N} D_n$ for some sequence of decidable subsets $D_n: S \to 2$. 
  Consider $U_n = \neg p(\neg D_n)$. We will prove that $C = \bigcap_{n:\N} U_n$.
  \item First we show that $C\subseteq U_n$ for any $n:\N$. If there is some $x:X$ such that $x\in C$ and $x \in q(\neg D_n)$, then by surjectivity of $q$, there is some $s:S$ with $s\in \neg D_n$ and $s\in q^{-1}(C)$, contradiciting $q^{-1}(C)= \bigcap_{n:\N} D_n$. Thus $C\subseteq \neg q(\neg D_n)$. 
  \item Now we show that $\bigcap_{n:\N} U_n \subseteq C$. 
    Let $x:X$ and suppose that for all $n:\N$ we have $x\in U_n$. Then for any $x:X$ and $s\in q^{-1}(x)$ we have for any $n:\N$ that $s \notin \neg D_n$, and by decidablity of $D$ hence $s\in D_n$. Thus $s\in q^{-1}(C)$ and $x\in C$ as required. 
\end{proof}

\begin{corollary}\label{approxDeltaByOpens}
  Let $X:\CHaus$, then there is some sequence of entourages $(U_n)_{n:\N}$ subsets of $X \times X$ such that $\bigcap_{n:\N} U_n = \Delta$. 
\end{corollary}

\begin{lemma}
  Extension for $\CHaus$. 
\end{lemma}
\begin{proof}
  By \Cref{approxDeltaByOpens}, there is a sequence of entourages $U_n\subseteq X \times X$ such that $\bigcap_{n:\N} U_n = \Delta$. 

  For each $n:\N$ we have $\Delta \subseteq U_n$, and hence by \Cref{TODO} there are opens $V_n$ and closeds $K_n$ such that $\Delta \subseteq V_n \subseteq K_n \subseteq U_n$. 

  By relative uniform continuity and dependent choice, we can pick for each $n:\N$ some $H_n$ entourage of $X$ such that for any $q':Q$ with $H_n(\iota q,\iota q')$ we have $V_n(f q , fq')$. 
  Define now for each $x:X$ the set $F_n(x) := \{y:Y | \forall_{q:Q} (H_n(\iota q , x) \to K_n(f(q),y)\}$. By \Cref{TODO}, this is a closed subset of $Y$. 
  We will now argue that $\bigcap_{n:\N} F_n(x)$ is a singleton, by showing it is a merely inhabited proposition. 
\item To see that it is inhabited, by \Cref{TODO} it is sufficient to show that $\bigcap_{n<k} F_n(x)$ is inhabited for all $k:\N$. So let $x:X$, consider $\{x':X |\forall_{n<k} H_n(x,x')\}$ this is an inhabited open subset, hence there exists some $q$ such that $H_n(x,q)$ for all $n:\N$. We claim that $f(q)\in \bigcap_{n<k} F_n(x)$. 
  Now suppose that $q':Q$ with $H_n(x, \iota q')$. As also $H_n(x , \iota q)$ we get \rednote{Assume H symmetric} $H_n^2(\iota q',\iota q)$, hence also $U(f q', fq)$ as required. 
\item To see that it is a proposition. Assume that for all $n:\N$ we have $y,y':F_n(x)$. By density, for each $n:\N$ we can find some $q:Q$ with $H_n(\iota q , x)$, hence for each $n:\N$ we have $K_n(f(q),y)$ and $K_n(f(q),y')$ and thus $K_n^2(y,y')$. 
  But then \rednote{Make it so $K_n^2\subseteq U_n$} we have $U_n(y,y')$ for all $n:\N$, hence $(y,y') \in \Delta$ and thus $y = y'$. 
\end{proof}
